% !PRINTED 419
\weberpage{419}
\begin{center}
\S.\,111.

\gesperrt{The Resolvents of the Icosahedral Equation.}
\end{center}

The formulation of the problem which we gave at the close of the preceding section leads us to the question of the various resolvents of the icosahedral equation. To every divisor of the icosahedral group there corresponds such a resolvent, whose degree is equal to the index of the divisor, and is therefore always a divisor of 60. Now the icosahedral group contains (\S.\,59):

\begin{tabular}{@{}ll c l@{}}
1) Tetrahedral groups (\S.\,56): & Resolvents & of the $5$th & degree \\
2) Dihedral groups $D_5$ (\S.\,55): & \textquotedbl & of the $6$th & \textquotedbl \\
3) Dihedral groups $D_3$: & \textquotedbl & of the $10$th & \textquotedbl \\
4) Cyclic groups $C_5$: & \textquotedbl & of the $12$th & \textquotedbl \\
5) Four-groups $D_2$: & \textquotedbl & of the $15$th & \textquotedbl \\
6) Cyclic groups $C_3$: & \textquotedbl & of the $20$th & \textquotedbl \\
7) Cyclic groups $C_2$: & \textquotedbl & of the $30$th & \textquotedbl \\
\end{tabular}

Since the icosahedral group is simple, all these equations are \gesperrt{total resolvents} of one another; we confine ourselves here to the consideration of the most important among them.

The resolvents of lowest, namely of degree $5$, are those belonging to the tetrahedral group.

To find them, we must seek a function of $\vartheta$ which remains unchanged by the substitutions of the tetrahedral group, while it is five-valued in the icosahedral group.

We take the icosahedral group in the normal form set up in \S.\,58, (22), and take out the tetrahedral group $Q$ (\S.\,59) contained in it, which we have constructed over the four-group
\weq{1}{1, \psi, \chi, \psi\chi}
in which, if $\varepsilon$ is an imaginary fifth root of unity, $\psi$ and $\chi$ have the meaning:
\weq{2}{\psi(x) = \dfrac{-1}{x}, \quad \chi(x) = \dfrac{\omega x+1}{x-\omega}, \quad \omega = \varepsilon + \varepsilon^{-1}.}

If we set $\Theta(x) = \varepsilon x$, the whole icosahedral group becomes
\weq{3}{P = Q + Q\Theta + Q\Theta^2 + Q\Theta^3 + Q\Theta^4,}
and
\weq{4}{Q_h = \Theta^{-h} Q \Theta^h}
are the tetrahedral groups conjugate to $Q$.
